\documentclass[tikz, border=6pt]{standalone}
% Shared preamble of the standalone TikZ figures (compiled by scripts/build_figures.py).
\usepackage{amsmath, amssymb}
\usepackage{xeCJK}
\setCJKmainfont[BoldFont={Hiragino Sans W6}]{Hiragino Sans W4}
\usetikzlibrary{arrows.meta, positioning, decorations.pathreplacing, calc, patterns}
\definecolor{rod}{RGB}{70,110,170}
\definecolor{rodb}{RGB}{200,120,50}
\definecolor{okc}{RGB}{60,140,80}
\definecolor{ngc}{RGB}{190,70,60}
\definecolor{boxc}{RGB}{235,240,248}
\tikzset{
  >={Stealth[length=2.4mm]},
  every node/.style={font=\small},
  box/.style={draw, rounded corners=2pt, fill=boxc, align=flush center, inner sep=4pt},
  rodA/.style={draw, fill=rod!35},
  rodB/.style={draw, fill=rodb!35},
  note/.style={font=\footnotesize, align=flush center},
  expl/.style={font=\small, align=flush left},
  brace/.style={decorate, decoration={brace, amplitude=4pt}},
}

\begin{document}
\begin{tikzpicture}
  \node[anchor=west, font=\bfseries] at (-0.2, 3.1) {(a) 単純系};
  \draw[box] (0,0.6) rectangle (3.2,2.6);
  \node[align=center] at (1.6,1.6) {$U, V, N$};
  \node[note] at (1.6, 0.15) {状態は一つに決まる};
  \begin{scope}[xshift=5cm]
    \node[anchor=west, font=\bfseries] at (-0.2, 3.1) {(b) 仕切りのある系};
    \foreach \k/\a/\b in {0/{$U/2$}/{$U/2$}, 1/{$U/3$}/{$2U/3$}} {
      \begin{scope}[xshift=\k*3.9cm]
        \draw[box] (0,0.6) rectangle (3.2,2.6);
        \draw[very thick] (1.6,0.6) -- (1.6,2.6);
        \node at (0.8,1.6) {\a};
        \node at (2.4,1.6) {\b};
      \end{scope}
    }
    \node[note, text width=7cm] at (3.55, 0.15) {全体の $U, V, N$ は同じでも、エネルギーの分け方の数だけ状態がある};
  \end{scope}
\end{tikzpicture}
\end{document}
